\section{Sobreexposición de personal o público}

\subsection*{Clasificación: II-A}

\subsubsection{Si la sobreexposición se sospecha porque uno de los dosímetros de lectura directa se encuentra saturado y no existe otro indicio de anormalidad}
Notifique al departamento de seguridad radiológica y proceda a leer el dosímetro TLD.

\subsubsection{Si los dosímetros tanto del auxiliar como del técnico se encuentran saturados o bien algunas radiografías han sido veladas, o se encuentra alguna anormalidad que indique la sobreexposición}
Proceda como sigue:

\begin{itemize}
    \item Asegúrese de que la fuente se encuentra segura en su contenedor.
    \item Haga una revisión mental de todas las maniobras ejecutadas y anote, especificando los siguientes datos:
    \begin{enumerate}
        \item Dosis que registraba el dosímetro al inicio de la jornada o la última vez que lo consultó.
        \item Número y tipo de maniobras.
        \item Tiempo aproximado de cada una de ellas.
        \item Distancia al contenedor y mangueras.
        \item Actividad de la fuente.
    \end{enumerate}
\end{itemize}